% =============================================================================
% quiz template — single in-class quiz with optional bank workflow
% =============================================================================
%
% Compile with:
%   lualatex template.tex
%
% This template demonstrates BOTH styles in one file:
%   1. Inline \question text (the original quiz pattern)
%   2. \getproblem retrieval from a problem bank
%
% Keep whichever you prefer; mix and match if useful.

\documentclass{quiz}

\meta{
	quiz-number = 1,

	% Optional: replace boxed instructions with a separate file, e.g.
	%   quiz-instructions-file = group-quiz-instructions
	% (filename without .tex; \input'd at \maketitle.)
}

% ---- Inline bank entries (or move into a separate `quiz_bank.tex` and
% ---- \loadbank{quiz_bank} below) -------------------------------------
\newproblem{eval-poly}{topic=eval, diff=easy}{
	Evaluate $3 \cdot (-2)^2 + 5 \cdot 4 - 3 \cdot 3^2$.
}[$5$]

\newproblem{linear-eq}{topic=algebra, diff=easy}{
	Solve $2x + 5 = 11$ for $x$.
}[$x = 3$]

% Alternatively, keep the bank in its own file:
% \loadbank{quiz_bank.tex}

\begin{document}
\maketitle

\begin{questions}

% --- (a) Bank-backed: pulled from \newproblem above -------------------
\question \getproblem{eval-poly}
	\begin{parts}
		\part[3] (Show your arithmetic.)
			\vspace{\stretch{1}}
		\part[2] State the order of operations rules you used.
			\vspace{\stretch{1}}
	\end{parts}

% --- (b) Bank-backed: random match from a key=value query -------------
\question \getproblem{topic=algebra, diff=easy}
	\vspace{\stretch{1}}

% --- (c) Inline (original quiz style — still works) -------------------
\question Solve the inequality $|x - 2| < 5$.
	\vspace{\stretch{1}}

\extracredit[2]{Find a real number $x$ such that $x^x = 1/4$.}
	\vspace{\stretch{1}}

\end{questions}

\end{document}
